\subsubsection{定理的陈述与证明}\label{subsec:ck-statement-proof}

考虑偏微分方程组
\begin{equation}\label{eq:ck-quasilinear-system}
\dfrac{\partial u_i}{\partial x_n}=\sum_{k=1}^{n-1}\sum_{j=1}^Na_{ij}^k(\bm{p})\dfrac{\partial u_j}{\partial x_k}+b_i(\bm{p}),\qquad i=1,\ldots,N,
\end{equation}
其中 $\bm{p}$ 表示向量 $(x_1,\ldots,x_{n-1},u_1,\ldots,u_N)$, 初始条件为
\begin{equation}\label{eq:ck-zero-initial-data}
u_i=0\quad\text{on }x_n=0,\qquad i=1,\ldots,N.
\end{equation}
将证明以下局部存在性结果.

\begin{Theorem}
\label{thm:cauchy-kovalevskaya-local-existence}
设函数 $a_{ij}^k$ 和 $b_i$ 在 $\mathbb{R}^{N+n-1}$ 的原点处实解析. 那么, 方程组 \eqref{eq:ck-quasilinear-system} 配以初始条件 \eqref{eq:ck-zero-initial-data}, 具有在原点处实解析的解组 $u_i$, 且该解组在实解析函数类中唯一.
\end{Theorem}

\begin{Proof}首先形式地计算 $u_i$ 在原点处的所有导数. 从 \eqref{eq:ck-zero-initial-data} 得知, 各阶切向导数全为零. 因而, 仅有的非零一阶导数沿 $x_n$ 方向, 且 \eqref{eq:ck-quasilinear-system} 给出 $\partial u_i/\partial x_n(\bm{0})=b_i(\bm{0})$. 接着对 \eqref{eq:ck-quasilinear-system} 求导, 可以得到 $u_i$ 的二阶导数. 得到
\begin{equation}\label{eq:ck-second-derivative-recurrence}
\begin{aligned}
\dfrac{\partial^2u_i}{\partial x_n\partial x_l}(\bm{0})
&=\sum_{k=1}^{n-1}\sum_{j=1}^Na_{ij}^k(\bm{0})\dfrac{\partial^2u_j}{\partial x_k\partial x_l}(\bm{0})\\
&\quad+\sum_{j=1}^{N+n-1}\dfrac{\partial b_i}{\partial p_j}(\bm{0})\dfrac{\partial p_j}{\partial x_l}(\bm{0}),
\end{aligned}
\end{equation}
由此, 当 $l=1,\ldots,n-1$ 时,
\begin{equation}\label{eq:ck-tangential-second-derivatives}
\dfrac{\partial^2u_i}{\partial x_n\partial x_l}=\dfrac{\partial b_i}{\partial p_l},
\end{equation}
当 $l=n$, $i=1,\ldots,N$ 时,
\begin{equation}\label{eq:ck-second-normal-derivative}
\begin{aligned}
\dfrac{\partial^2u_i}{\partial x_n^2}
&=\sum_{k=1}^{n-1}\sum_{j=1}^Na_{ij}^k(\bm{0})\dfrac{\partial^2u_j}{\partial x_k\partial x_n}(\bm{0})
+\sum_{j=1}^N\dfrac{\partial b_i}{\partial p_{n-1+j}}(\bm{0})\dfrac{\partial u_j}{\partial x_n}(\bm{0})\\
&=\sum_{k=1}^{n-1}\sum_{j=1}^Na_{ij}^k(\bm{0})\dfrac{\partial b_j}{\partial p_k}(\bm{0})
+\sum_{j=1}^N\dfrac{\partial b_i}{\partial p_{n-1+j}}(\bm{0})b_j(\bm{0}).
\end{aligned}
\end{equation}


类似地进行下去, 可以计算各阶导数. $u_i$ 在原点处的任意导数的表达式, 都是由 $a_{ij}^k$ 和 $b_i$ 在原点处的导数组成的正系数多项式.

于是可以构造 $u_i$ 的形式 Taylor 级数. 为证明这一 Taylor 级数收敛, 使用优函数法. 令
\begin{equation}\label{eq:ck-coefficient-majorants}
a_{ij}^k(\bm{p})\preceq A_{ij}^k(\bm{p}),\qquad b_i(\bm{p})\preceq B_i(\bm{p}),
\end{equation}
并令 $U_i$ 为如下问题的解:
\begin{equation}\label{eq:ck-majorant-system}
\dfrac{\partial U_i}{\partial x_n}=\sum_{k=1}^{n-1}\sum_{j=1}^NA_{ij}^k(\bm{p})\dfrac{\partial U_j}{\partial x_k}+B_i(\bm{p}),\qquad i=1,\ldots,N,
\end{equation}
初始条件为
\begin{equation}\label{eq:ck-majorant-zero-data}
U_i=0\quad\text{on }x_n=0,\qquad i=1,\ldots,N.
\end{equation}
那么显然对任意 $\alpha$ 有 $|D^\alpha u_i(\bm{0})|\leq D^\alpha U_i(\bm{0})$; 因而, 如果 $U_i$ 解析, $u_i$ 的形式幂级数就收敛.

还需构造适当的函数 $A_{ij}^k$ 和 $B_i$. 假设 $a_{ij}^k$ 和 $b_i$ 属于 $C_{M,r}(\bm{0})$. 那么可以由 \eqref{eq:analytic-rational-majorant} 找到一个优函数, 即
\begin{equation}\label{eq:ck-rational-coefficient-majorants}
A_{ij}^k=B_i=\dfrac{Mr}{r-x_1-\cdots-x_{n-1}-U_1-\cdots-U_N}.
\end{equation}
也就是说, 需要考虑初值问题
\begin{equation}\label{eq:ck-rational-majorant-system}
\dfrac{\partial U_i}{\partial x_n}=\dfrac{Mr}{r-x_1-\cdots-x_{n-1}-U_1-\cdots-U_N}\left(1+\sum_{k=1}^{n-1}\sum_{j=1}^N\dfrac{\partial U_j}{\partial x_k}\right),
\end{equation}
及初始条件 \eqref{eq:ck-majorant-zero-data}.

可以寻找形如 $U_i(\bm{x})=V(x_1+\cdots+x_{n-1},x_n)$ 的解; 由此得到 $V(s,t)$ 的方程为
\begin{equation}\label{eq:ck-scalar-majorant-problem}
V_t=\dfrac{Mr}{r-s-NV}(1+N(n-1)V_s),\qquad V(s,0)=0.
\end{equation}


此方程可显式求解. 首先注意, 特征线 $s=s(t)$ 由方程
\begin{equation}\label{eq:ck-majorant-characteristics}
\dfrac{\mathrm{d}s}{\mathrm{d}t}=-\dfrac{(n-1)NMr}{r-s-NV(s(t),t)}.
\end{equation}
给出.

沿一条特征线, 有
\begin{equation}\label{eq:ck-majorant-along-characteristic}
\dfrac{\mathrm{d}}{\mathrm{d}t}V(s(t),t)=\dfrac{Mr}{r-s(t)-NV(s(t),t)}=-\dfrac1{(n-1)N}\dfrac{\mathrm{d}s}{\mathrm{d}t}.
\end{equation}
对 $t$ 积分, 得到
\begin{equation}\label{eq:ck-majorant-characteristic-integral}
V(s(t),t)=-\dfrac1{(n-1)N}s(t)+\dfrac1{(n-1)N}s(0).
\end{equation}


将 \eqref{eq:ck-majorant-characteristic-integral} 代入 \eqref{eq:ck-majorant-characteristics}, 得到
\begin{equation}\label{eq:ck-reduced-characteristic-equation}
\dfrac{\mathrm{d}s}{\mathrm{d}t}=-\dfrac{(n-1)MNr}{r-s+\dfrac{s}{n-1}-\dfrac{s(0)}{n-1}},
\end{equation}
积分可得
\begin{equation}\label{eq:ck-initial-characteristic-coordinate}
s(0)=\dfrac{n-1}{n}\left[r+\dfrac1{n-1}s-\sqrt{(r-s)^2-2nNMrt}\right].
\end{equation}
将它代入 \eqref{eq:ck-majorant-characteristic-integral}, 最终得到
\begin{equation}\label{eq:ck-explicit-scalar-majorant}
V(s,t)=\dfrac1{Nn}\left(r-s-\sqrt{(r-s)^2-2nNMrt}\right).
\end{equation}
这个表达式在 $s=t=0$ 处解析. 证明完毕.\end{Proof}

这一定理保证, 初始曲面的每一点的某个邻域内都存在解. 取所有这些邻域的并, 就得到在包含初始曲面的一个开集内存在解. 还要注意, 解的 Taylor 级数保证在某个球内收敛, 球的半径只依赖于 $n$, $N$, $M$ 和 $r$.

强调, 这一定理具有严格的局部性质, 即解可能在 $x_n$ 的某个有限值处不再存在, 或至少失去解析性. 后面研究激波形成时, 将看到这方面的例子. 此外, 即使在初始曲面的一个邻域内, 这一定理也不排除非解析解的存在. 唯一性只在解析函数类中得到保证. 下一节将证明的 Holmgren 定理断言, 对线性方程, 解析解在一个更大的函数类中也是唯一的.
